\chapter*{Preface}
\addcontentsline{toc}{chapter}{Preface}

Meaning, across the three books that precede this one, has been named repeatedly without being pinned down. Book 1 identified it, in passing, with a closed but non-exact cochain; Book 3 showed how graded asymmetry seeds the directional structure without which nothing could persist long enough to be called meaningful at all. The present book takes up the identification Book 1 only gestured toward and develops it as the central object of a full cohomology theory. Meaning, on the account defended here, is the only curvature the plenum conserves absolutely: potential dissipates, entropy accumulates, but a non-trivial cohomology class, once established, cannot be destroyed by any process the RSVP framework recognizes as smoothing. It can only be forgotten, in a precise technical sense developed in Part V, and forgetting, on this account, is a different and more contingent fact than destruction.

\chapter*{Concept Map}
\addcontentsline{toc}{chapter}{Concept Map}

Three notions organize this book: entropy, already familiar from Books 1 and 3 as the field whose gradients drive curvature and whose global increase is never violated; cohomology, introduced here as the precise sense in which some configurations of the fields survive that increase while others do not; and ethical invariance, the claim, defended across Part III, that moral coherence is itself a cohomological fact rather than a separate normative overlay on the physics. The chapters below move from the first of these to the second and from the second to the third, and the closing parts show that cognition and culture are best understood as domains in which the same cohomological machinery is already operating, whether or not it is recognized as such.

\chapter*{Reading Note}
\addcontentsline{toc}{chapter}{Reading Note}

Mathematical rigor and interpretive reflection are co-equal modes of proof in this book, not sequential stages in which rigor comes first and reflection is added afterward as commentary. Each chapter accordingly closes with a short passage, marked as a Meaning Note, stating in condensed form what the chapter's formal result shows to persist through smoothing. These notes are not summaries; they are the point at which the chapter's mathematics and its ethical content are shown to coincide.

\part{The Cohomological Plenum}

\chapter{Closed Forms and Semantic Persistence}

The RSVP differential complex is the sequence
\[
0 \longrightarrow \PhiField \xrightarrow{d_1} \vField \xrightarrow{d_2} \SEntropy \longrightarrow 0,
\]
with $d_2 \circ d_1 = 0$, so that the composite map from potential to entropy, passing through flow, vanishes identically. This nilpotency is not an additional assumption but a restatement, in cochain language, of the same Jacobi identity that governed the commutator algebra of Book 1; the present chapter's task is to show that this restatement carries genuine content once the complex is used to classify which configurations of $\RSVPTriple$ survive the plenum's entropic evolution and which do not.

The physical interpretation of the complex is conservation of flux across the RSVP triple: whatever passes from potential through flow arrives at entropy without residue, and it is exactly this residue-free passage that licenses treating $d_1$ and $d_2$ as a genuine differential rather than as a bookkeeping device. A closed form, in this complex, is a stable pattern, a configuration that $d_2$ or $d_1$ leaves unchanged in the sense relevant to cohomology; an exact form is a transient signal, one that arises only as the image of the preceding map and therefore carries no content beyond what that preceding stage already supplied. The cognitive metaphor proposed here, developed at length in Part IV, treats understanding as a closed trajectory in information space: a state of understanding, once reached, is not eroded by further processing in the way a merely derivative, exact signal is, and this difference in survival is precisely what the cohomological classification of this chapter is built to track.

\emph{Meaning Note.} What survives this chapter's construction is not any particular configuration of the fields but the distinction between closed and exact forms itself: a pattern's survival through smoothing is decided by which side of that distinction it falls on, not by any further property that must be checked case by case.

